%!TEX root = ../main.tex

% ============================================================
%  CAPÍTULO 8: CONCLUSIÓN
% ============================================================

\chapter{Conclusiones}\label{ch:conclusion}

% ──────────────────────────────────────────────────────────────
\section{Respuesta a la pregunta de investigación}
% ──────────────────────────────────────────────────────────────

Este documento se originó a partir de una inquietud central: evaluar si era metodológicamente posible diseñar un modelo de \textit{nowcasting} electoral elección-agnóstico y emancipado de las encuestas para el complejo entorno local colombiano. A la luz de los hallazgos empíricos, la respuesta global es afirmativa con matices estructurales: el modelo es conceptual y operativamente posible, pero su techo predictivo actual está condicionado por el volumen histórico de entrenamiento.

En primer lugar, se demostró irrefutablemente que las interacciones ciudadanas en plataformas sociales \textbf{tienen un poder predictivo real} sobre el comportamiento del electorado colombiano. Las pruebas no paramétricas arrojaron diferencias significativas ($p < 0.05$) entre ganadores y perdedores, y la correlación de rango cruzó un umbral robusto de Spearman $\rho > 0.46$. No obstante, la experimentación también comprobó que, frente a un problema de alta varianza como una elección fragmentada, la señal extraída aisladamente de la red sin covariables de apoyo es ruidosa e insuficiente para dictar sentencias determinísticas absolutas a nivel nacional.

En segundo término, la viabilidad del modelo descansaba sobre la premisa de la \textbf{reconstrucción temporal}. Las pruebas sobre el panel de control validaron que la interpolación matemática mediante distribuciones de Weibull es excepcionalmente precisa (inyectando errores menores al 20\% en más del 89\% del material muestral observado tempranamente). En consecuencia, un marco analítico elección-agnóstico se ha consolidado como un enfoque empírico viable y robusto: el modelo puede abstraer y aplicar reglas universales. Sin embargo, para que las arquitecturas topológicas (como redes neuronales o ensambles complejos) logren trascender el desempeño del algoritmo \textit{baseline} trivial basado en dominancia relativa, se torna imperativo inyectar iteraciones masivas de eventos históricos en la matriz de aprendizaje. 

% ──────────────────────────────────────────────────────────────
\section{Lecciones aprendidas}
% ──────────────────────────────────────────────────────────────

A lo largo del desarrollo computacional y la revisión estadística del proyecto, se decantaron lecciones críticas para la práctica del \textit{nowcasting} digital en el sur global:

\begin{itemize}
    \item \textbf{La supremacía de la dominancia relativa:} Quedó evidenciado que el ecosistema de campaña funciona como un juego de suma cero territorial. La posición de un candidato relativa a la conversación local de su ciudad (\textit{share of voice}) es una métrica inmensamente superior al volumen absoluto de interacciones, pues este último es susceptible de distorsiones poblacionales y mediáticas inherentes a las dinámicas atípicas de los perfiles públicos masivos.
    \item \textbf{El paradigma de clasificación:} Frente a matrices limitadas por la escasez de datos subnacionales etiquetados, plantear la predicción como un problema de clasificación probabilística discreta (\textit{¿quién se lleva la elección?}) probó ser cualitativamente más realista y convergente que pretender establecer márgenes porcentuales directos (\textit{¿cuánto saca?}) a través de regímenes de regresión.
    \item \textbf{El abismo de la última semana:} El obstáculo empírico más resistente para las métricas agregativas continúa siendo el giro coyuntural tardío. La explosión de interacciones preexistentes oculta frecuentemente las volatilidades sutiles de la opinión pública suscitadas por eventos, escándalos o alianzas que acontecen en los tres días anteriores al evento de urnas.
    \item \textbf{Autonomía computacional:} La comprobación de un ciclo de vida completo de análisis —desde la orquestación y el \textit{scraping} (\textit{Make.com}, \textit{Playwright}, API de X, \textit{Apify}) hasta la proyección asintótica y el modelado matricial— probó que este \textit{pipeline} analítico complejo es enteramente ejecutable desde laboratorios comerciales e independientes que operen con un coste logístico nominal o nulo.
\end{itemize}

% ──────────────────────────────────────────────────────────────
\section{Próximos pasos}
% ──────────────────────────────────────────────────────────────

La natural prolongación de este cuerpo de investigación consiste en llevar la experimentación a la fase de masificación de datos, integrando \textit{datasets} de contiendas asimétricas de más naciones de América Latina y diferentes marcos temporales, apuntalando así el tamaño de las muestras para destrabar el uso de aprendizaje profundo. Inminentemente, se hace obligatoria la recalibración inter-regional del modelo Weibull, así como el rediseño dimensional para inyectar vectores categóricos de sentimiento. La culminación de este esfuerzo empírico se cristalizará durante la apertura de las urnas de las próximas elecciones territoriales en Colombia de 2027, las cuales fungirán como el juez definitivo ciego (\textit{out-of-time sample}) sobre el cual se medirá, bajo estricto rigor metodológico ex-ante, la madurez proyectiva de la inteligencia desarrollada.
